\section{Semantic Category}
\label{sec:category}

\subsection{Free Category and Universal Property}

\begin{definition}
$\mathcal{A}_1(\Sigma) = \operatorname{Free}(\Delta_0, \Delta_1, d_0, d_1)$. The free category generated by the $0$-simplices (objects) and $1$-simplices (generating morphisms).
\end{definition}

\begin{theorem}
Category axioms + universal property. \textup{[Lean~4]}
\end{theorem}

\begin{remark}
The universal property is the theoretical foundation of the plugin architecture: any structure-preserving map from generators extends uniquely to a functor.
\end{remark}

\subsection{2-Category and Refactoring Safety}

\begin{definition}
$\mathcal{A}_2(\Sigma)$. 2-morphisms = equivalences between two paths (two dependency chains).
\end{definition}

\begin{remark}
A 2-morphism exists $\iff$ replacement is safe. For $k \geq 3$, connections to homotopy type theory exist, but $k = 2$ suffices for knowledge management.
\end{remark}

\subsection{Nerve, Induced Complex, and Structural Gap}

\begin{definition}
$\Delta^{\mathrm{ind}} = N(\mathcal{A}_1)$: the nerve of the free category (induced complex). $\Delta^{\mathrm{raw}} = \Delta(\Sigma)$: the original complex from references.
\end{definition}

\begin{definition}
$\Delta^+$ (implicit higher-order simplices), $\Delta^-$ (irreducible higher-order simplices), $\epsilon$ (orientation consistency error).
\end{definition}

\subsection{Duality Sharp/Flat}

\begin{definition}
$X^\sharp$: follow references. $X^\flat$: follow reverse references. Dual operators pervade all tools.
\end{definition}

\begin{theorem}
$\chi(\Delta) = \chi(\Delta^{\mathrm{op}})$.
\end{theorem}

\subsection{Enrichment}

\begin{definition}[Enrichment]
\label{def:enrichment}
An \emph{enrichment} is a function $E : R \to R$ that modifies only the auxiliary part:
\[
E(r)|_{K_{\mathrm{ess}}} = r|_{K_{\mathrm{ess}}}, \qquad h_{\mathrm{id}}(E(r)) = h_{\mathrm{id}}(r).
\]
Enrichments compose: $E_2 \circ E_1$ is again an enrichment, and enrichments writing to disjoint keys commute ($E_1 \circ E_2 = E_2 \circ E_1$). Each enrichment is an endofunctor on the semantic category that acts as the identity on objects and morphisms while extending the key-value content.
\end{definition}

\subsection{Maps and Import}

\begin{definition}[Semantic map]
A reference-preserving map between semantic structures.
\end{definition}

\begin{theorem}
Preserves composition and faithfulness. \textup{[Lean~4]}
\end{theorem}

\begin{theorem}
Faithful import preserves homology.
\end{theorem}

\subsection{Quotient, Core, and Hierarchy Tower}

\begin{definition}
Quotient (arbitrary dimension). Core structure (structural fingerprint).
\end{definition}

\begin{theorem}
Quotients do not increase homological dimension.
\end{theorem}

\begin{corollary}
The hierarchy tower terminates finitely; $k^*$ decreases monotonically at each level.
\end{corollary}

\begin{corollary}
The embedding dimension of the core $\leq$ the embedding dimension of the original.
\end{corollary}

\begin{remark}
Time series of cores = evolutionary history. A jump in the core = structural transformation.
\end{remark}

\subsection{Slices, Localization, and Sheaves}

\begin{definition}
$\operatorname{St}^\sharp(a)$, $\operatorname{St}^\flat(a)$: full-dimensional dependency tree / influence range.
\end{definition}

\begin{definition}[Localization]
$\mathcal{S}[W^{-1}]$: inverting equivalences, versions, or organizational structure. Used with the snake lemma.
\end{definition}

\begin{definition}[Sheaf condition]
Grothendieck topology + sheaf condition = type checking.
\end{definition}

\subsection{Functionals}

\begin{definition}
$k$-functional $f : \Delta_k \to \mathbb{R}$. Unified with cohomology: $C^k$ is the space of $k$-functionals, $H^k$ is the space of essentially distinct $k$-functionals.
\end{definition}

\begin{definition}
$F^\sharp$, $F^\flat$: push-forward and pull-back of functionals.
\end{definition}
